% =========================================================
% Proof Stub: Cauchy is a Tail Property
% Source: volume-iii/analysis/sequences/notes/sequences/notes-k-tails.tex
% Mode: standard
% =========================================================

\clearpage
\phantomsection
\hypertarget{proof-RA-ABB-T-cauchy-tail-prop}{}

\subsubsection[Cauchy is a Tail Property]{Proof --- Cauchy is a Tail Property}

\bigskip

\noindent
\textbf{Source.}
See \texttt{notes-k-tails.tex}.

\vspace{0.75em}

\noindent
\textbf{Goal.}
$(x_n)$ is Cauchy if and only if every $k$-tail is Cauchy.

\vspace{0.75em}

\noindent
\textbf{Logical form.}
\[
  (x_n) \text{ Cauchy} \Longleftrightarrow \forall k,\; (x_n)_{n \ge k} \text{ Cauchy}.
\]

\vspace{0.75em}

\noindent
\textbf{Proof strategy.}
(⇒) Any $k$-tail is a subsequence of a Cauchy sequence, hence Cauchy. (⇐) The $0$-tail is the whole sequence; if it is Cauchy, the sequence is Cauchy.

\vspace{0.75em}

\noindent
\textbf{Proof.}
\begin{proof}
\Claim $(x_n)$ is Cauchy if and only if every $k$-tail is Cauchy.

\Given 

\Goal 

\Strategy (⇒) Any $k$-tail is a subsequence of a Cauchy sequence, hence Cauchy. (⇐) The $0$-tail is the whole sequence; if it is Cauchy, the sequence is Cauchy.

\medskip

% --- (⇒) given ε > 0 get N from Cauchy condition; for m,n ≥ max(N,k): |x_m - x_n| < ε ---
% --- (⇐) the full sequence (x_n) is its own 0-tail ---

\end{proof}

\begin{remark*}[Proof shape]
The Cauchy condition depends only on the tail (indices $\ge N$), so shifting the start index changes nothing.
\end{remark*}

\begin{remark*}[Dependencies]
\begin{itemize}
  \item Definition of Cauchy sequence.
  \item Cauchy sequences are inherited by subsequences (since any tail is a subsequence).
\end{itemize}
\end{remark*}
